\documentclass[11pt]{article}

% ---------- geometry ----------
\usepackage[margin=1in]{geometry}

% ---------- math ----------
\usepackage{amsmath,amssymb,amsthm,mathtools}

% ---------- figures / tables ----------
\usepackage{graphicx}
\usepackage{booktabs}
\usepackage{subcaption}

% ---------- algorithms ----------
\usepackage[ruled,vlined,linesnumbered]{algorithm2e}

% ---------- typography ----------
\usepackage{microtype}
\usepackage[T1]{fontenc}
\usepackage{lmodern}

% ---------- references ----------
\usepackage[numbers,sort&compress]{natbib}
\usepackage[colorlinks,citecolor=blue,linkcolor=blue,urlcolor=blue]{hyperref}
\usepackage{cleveref}

% ---------- misc ----------
\usepackage{xcolor}
\usepackage{enumitem}

% ---------- custom commands ----------
\newcommand{\R}{\mathbb{R}}
\newcommand{\E}{\mathbb{E}}
\newcommand{\KL}{\mathrm{KL}}
\newcommand{\clip}{\mathrm{clip}}
\DeclareMathOperator{\softmax}{softmax}
\newcommand{\todo}[1]{{\color{red}\textbf{[TODO: #1]}}}

% ==========================================================================
\title{Holistic vs.\ Per-Step PPO for Stochastic Flow Policies\\in Continuous Control}
\author{
  Jason Trinh \\
  University of California, Berkeley \\
  \texttt{jason.trinh@berkeley.edu}
}
\date{Spring 2026 --- CS 185 Final Project}

\begin{document}
\maketitle

\begin{abstract}
Generative flow policies parameterize actions as the output of a multi-step
stochastic process, where noise is iteratively refined into an action through a
learned velocity field.  Proximal Policy Optimization (PPO) typically treats this
chain as a black box, computing a single importance-sampling ratio from the
holistic action density.  We ask whether PPO can instead operate \emph{per step}
of the generative chain, and whether non-uniform credit assignment across steps
improves learning.  We formalize per-step PPO for stochastic flow policies,
propose uniform, learned-global, and state-dependent weighting schemes, and
compare all variants on continuous-control benchmarks (HalfCheetah-v5,
Hopper-v5).  Our experiments show \todo{one-sentence summary of main result}.
\end{abstract}

% ---------- sections ----------
\section{Introduction}
\label{sec:intro}

Generative models have emerged as a powerful policy class for continuous
control.  Diffusion policies~\citep{chi2023diffusion} and flow-matching
policies~\citep{lipman2023flow} parameterize the action distribution as the
output of an iterative denoising or transport process, enabling expressive
multimodal action distributions that outperform simple Gaussian policies on
complex manipulation and locomotion tasks.

When such generative policies are fine-tuned with on-policy reinforcement
learning---specifically Proximal Policy Optimization
(PPO)~\citep{schulman2017ppo}---the standard practice is to treat the entire
generative chain as a \emph{black box}.  The holistic log-probability
$\log p_\theta(\mathbf{a}\mid\mathbf{s})$ is obtained by summing per-step
log-conditionals and a change-of-variable correction, and a single
importance-sampling ratio $r_t = p_\theta(\mathbf{a}_t\mid\mathbf{s}_t) /
p_{\theta_{\text{old}}}(\mathbf{a}_t\mid\mathbf{s}_t)$ drives the PPO
surrogate loss.

However, a stochastic flow policy with $K$ internal steps exposes $K$
\emph{tractable} conditional densities
$p_\theta(\mathbf{z}_{k-1}\mid\mathbf{z}_k,\mathbf{s},k)$, each a simple
Gaussian.  This raises a natural design question:

\begin{quote}
\emph{Should PPO operate holistically on the final action, or per-step on each
internal transition?  And if per-step, should every step receive equal credit
for the task reward, or should credit be weighted?}
\end{quote}

Per-step optimization is appealing because it decomposes a single, potentially
high-variance ratio into $K$ smaller, better-behaved ratios, each of which can
be individually clipped.  Weighted credit assignment further allows the
algorithm to focus gradient signal on the steps that matter most for downstream
performance.

\paragraph{Contributions.}
We make the following contributions:
\begin{enumerate}[leftmargin=*,nosep]
  \item We \textbf{formalize per-step PPO} for stochastic flow policies,
        deriving per-step importance-sampling ratios and clipped surrogate
        objectives (\Cref{sec:method}).
  \item We \textbf{propose weighted credit-assignment} schemes---uniform,
        learned-global, and state-dependent---that modulate how the task
        advantage is distributed across generative steps.
  \item We provide an \textbf{empirical comparison} on continuous-control
        benchmarks (HalfCheetah-v5, Hopper-v5), analyzing learning speed,
        final performance, clip fractions, and learned weight profiles.
\end{enumerate}

\section{Related Work}
\label{sec:related}

\paragraph{Proximal Policy Optimization.}
PPO~\citep{schulman2017ppo} is the workhorse on-policy algorithm for deep RL,
using clipped importance-sampling ratios to constrain policy updates.  Earlier
trust-region methods include TRPO~\citep{schulman2015trpo}.  Our work retains
the PPO framework but changes the granularity at which ratios and clipping are
applied.

\paragraph{Generative Models as Policies.}
Diffusion models~\citep{ho2020ddpm,song2021scorebased} have been adopted as
expressive policy classes.  \citet{chi2023diffusion} introduced Diffusion
Policy for imitation learning in robotics.  \citet{black2024ddpo} proposed DDPO,
which fine-tunes diffusion models with policy-gradient methods by treating each
denoising step as an RL action.  Flow-matching
policies~\citep{lipman2023flow,liu2023flow} replace the diffusion SDE with a
simpler ODE-based transport, offering faster sampling and simpler training
objectives.

\paragraph{Stochastic Interpolants and Stochastic Flow Matching.}
\citet{albergo2023stochastic} introduced stochastic interpolants, providing a
unified framework that generalizes both diffusion and flow matching by
interpolating between source and target distributions with optional stochasticity.
\citet{albergo2023building} extended this to build normalizing flows with
stochastic dynamics.  Our stochastic flow policy draws on this framework,
retaining per-step stochasticity so that each conditional
$p_\theta(\mathbf{z}_{k-1}\mid\mathbf{z}_k)$ has tractable Gaussian density.

\paragraph{RL Fine-Tuning of Generative Policies.}
Several works have applied RL to fine-tune generative policies.
\citet{black2024ddpo} compute per-step importance weights within DDPO.
\citet{fan2024dpok} use KL-regularized RL objectives for diffusion fine-tuning.
\citet{clark2024directly} study direct reward fine-tuning of diffusion models.
Most recently, DPPO~\citep{ren2025dppo} formulates diffusion policy optimization
as a two-level MDP and applies PPO with per-step structure, while
ReinFlow~\citep{hu2025reinflow} fine-tunes flow-matching policies with PPO using
injected stochasticity.  Our work is complementary: rather than proposing a new
algorithm, we conduct a controlled comparison of \emph{how} PPO credit and
clipping should be distributed across internal steps, testing uniform, learned
global, and state-dependent weighting schemes that these prior works do not
systematically evaluate.

\paragraph{Credit Assignment in RL.}
Temporal credit assignment is a classical challenge in
RL~\citep{sutton2018rl,minsky1961steps}.  Hindsight credit assignment
methods~\citep{harutyunyan2019hindsight,hung2019optimizing} attempt to
redistribute reward to relevant earlier decisions.  Our per-step weighting can
be viewed as a form of \emph{intra-action} credit assignment: distributing the
task-level advantage across the internal generative steps that produced the
action.

\paragraph{Positioning.}
Our work bridges generative-model fine-tuning and PPO implementation design.
DPPO~\citep{ren2025dppo} and ReinFlow~\citep{hu2025reinflow} demonstrate that
per-step PPO is viable for generative policies, but neither systematically
compares credit-assignment weighting schemes.  Our contribution is this
controlled comparison: we hold the policy architecture and PPO framework fixed
and vary only whether clipping and credit are applied holistically, per-step
uniformly, or per-step with learned (global or state-dependent) weights.

\section{Method}
\label{sec:method}

We consider on-policy reinforcement learning with a stochastic flow policy
$\pi_\theta$.  At each environment step, the agent observes state
$\mathbf{s}\in\R^{d_s}$ and produces action
$\mathbf{a}\in[-1,1]^{d_a}$ by running a $K$-step stochastic generative
chain.

% ------------------------------------------------------------------
\subsection{Stochastic Flow Policy}
\label{sec:method:flow}

The generative chain proceeds from pure noise to an action as follows.
Let $\Delta t = 1/K$ and let $\sigma > 0$ be the per-step noise scale.

\begin{enumerate}[leftmargin=*,nosep]
  \item Sample initial noise: $\mathbf{z}_K \sim \mathcal{N}(\mathbf{0},\mathbf{I})$.
  \item For $k = K, K{-}1, \ldots, 1$, compute the next latent:
\end{enumerate}
\begin{align}
  \mathbf{z}_{k-1}
    &= \mathbf{z}_k
       + \Delta t \cdot v_\theta(\mathbf{s}, \mathbf{z}_k, k)
       + \sigma\,\boldsymbol{\epsilon}_k,
  \qquad \boldsymbol{\epsilon}_k \sim \mathcal{N}(\mathbf{0},\mathbf{I}).
  \label{eq:flow_step}
\end{align}
\begin{enumerate}[leftmargin=*,nosep,resume]
  \item Apply a squashing function to obtain the action:
        $\mathbf{a} = \tanh(\mathbf{z}_0)$.
\end{enumerate}

Each transition has a tractable Gaussian conditional:
\begin{align}
  p_\theta(\mathbf{z}_{k-1} \mid \mathbf{z}_k, \mathbf{s}, k)
    &= \mathcal{N}\!\bigl(
         \mathbf{z}_{k-1};\;
         \mathbf{z}_k + \Delta t\, v_\theta(\mathbf{s}, \mathbf{z}_k, k),\;
         \sigma^2 \mathbf{I}
       \bigr).
  \label{eq:per_step_density}
\end{align}

The holistic log-density of the final action is:
\begin{align}
  \log p_\theta(\mathbf{a}\mid\mathbf{s})
    &= \sum_{k=1}^{K}
       \log p_\theta(\mathbf{z}_{k-1}\mid\mathbf{z}_k,\mathbf{s},k)
       + \log\!\left|\det\frac{\partial\,\tanh(\mathbf{z}_0)}
                               {\partial\,\mathbf{z}_0}\right|.
  \label{eq:holistic_logprob}
\end{align}

% ------------------------------------------------------------------
\subsection{Holistic PPO (Baseline)}
\label{sec:method:holistic}

Standard PPO computes a single importance-sampling ratio per environment
transition $(s_t, a_t)$:
\begin{align}
  r_t
    &= \frac{p_\theta(\mathbf{a}_t\mid\mathbf{s}_t)}
            {p_{\theta_{\text{old}}}(\mathbf{a}_t\mid\mathbf{s}_t)},
  \label{eq:holistic_ratio}
\end{align}
and optimizes the clipped surrogate:
\begin{align}
  L^{\text{CLIP}}(\theta)
    &= \E_t\Bigl[
         \min\!\bigl(
           r_t\, \hat{A}_t,\;
           \clip(r_t, 1{-}\epsilon, 1{+}\epsilon)\,\hat{A}_t
         \bigr)
       \Bigr],
  \label{eq:ppo_clip}
\end{align}
where $\hat{A}_t$ is the GAE advantage estimate and $\epsilon$ is the clip
parameter (typically $0.2$).

Because $r_t$ is the product of $K$ per-step ratios (plus a correction),
it can be very large or very small, making clipping less effective for
large $K$.

% ------------------------------------------------------------------
\subsection{Per-Step PPO}
\label{sec:method:perstep}

We propose to apply PPO's clipped surrogate \emph{independently at each
generative step}.  Define the per-step ratio:
\begin{align}
  r_{t,k}
    &= \frac{p_\theta(\mathbf{z}_{t,k-1}\mid\mathbf{z}_{t,k},\mathbf{s}_t,k)}
            {p_{\theta_{\text{old}}}(\mathbf{z}_{t,k-1}\mid\mathbf{z}_{t,k},\mathbf{s}_t,k)}.
  \label{eq:perstep_ratio}
\end{align}

Each step $k$ is assigned an advantage $A_{t,k} = w_k \cdot \hat{A}_t$, where
$w_k$ is a credit-assignment weight satisfying $\sum_k w_k = 1$.  The per-step
clipped surrogate is:
\begin{align}
  L^{\text{PS-CLIP}}(\theta)
    &= \E_t\!\left[
         \sum_{k=1}^{K}
         \min\!\bigl(
           r_{t,k}\, A_{t,k},\;
           \clip(r_{t,k}, 1{-}\epsilon, 1{+}\epsilon)\, A_{t,k}
         \bigr)
       \right].
  \label{eq:perstep_clip}
\end{align}

By clipping each $r_{t,k}$ individually, we prevent any single step from
producing an excessively large update, even when the holistic ratio $r_t$
would be far from~$1$.

% ------------------------------------------------------------------
\subsection{Weighted Credit Assignment}
\label{sec:method:weights}

We study three weighting schemes for the per-step advantages
$A_{t,k} = w_k \cdot \hat{A}_t$:

\paragraph{Uniform.}
All steps receive equal credit:
\begin{align}
  w_k &= \frac{1}{K}, \qquad \forall\, k.
  \label{eq:uniform}
\end{align}

\paragraph{Learned Global.}
A single set of learnable logits $\boldsymbol{\alpha}\in\R^K$ produces
step-dependent but state-independent weights:
\begin{align}
  w_k &= \softmax(\boldsymbol{\alpha})_k.
  \label{eq:learned_global}
\end{align}
The logits $\boldsymbol{\alpha}$ are trained end-to-end alongside $\theta$.

\paragraph{State-Dependent.}
A small network $h_\phi$ takes the current state and latent as input and
produces per-step weights that depend on the rollout:
\begin{align}
  w_{t,k} &= \softmax\!\bigl(
                h_\phi(\mathbf{s}_t, \mathbf{z}_{t,k}, k)
              \bigr)_k.
  \label{eq:state_dep}
\end{align}
Here $\phi$ denotes the parameters of the weighting network, which are
jointly optimized with $\theta$.

% ------------------------------------------------------------------
\subsection{Intra-Chain Temporal Credit Assignment}
\label{sec:method:intra}

The weighting schemes above split the environment advantage $\hat{A}_t$
into per-step shares using fixed or learned proportions, but these shares
are heuristic---they do not reflect which steps \emph{actually improved}
the latent toward a good action.

We propose \textbf{intra-chain temporal credit assignment}: we learn a
value function $V_\psi(\mathbf{s}, \mathbf{z}_k, k)$ that estimates the
expected action quality at each internal step, and compute genuine
per-step advantages via temporal differencing within the chain:
\begin{align}
  A_{t,k}^{\text{intra}}
    &= V_\psi(\mathbf{s}_t, \mathbf{z}_{t,k-1}, k{-}1)
     - V_\psi(\mathbf{s}_t, \mathbf{z}_{t,k}, k).
  \label{eq:intra_adv}
\end{align}
Intuitively, $A_{t,k}^{\text{intra}}$ measures how much step $k$
improved the latent---a step that moves the latent toward a
high-value region of the chain receives positive credit, while a step
that degrades quality receives negative credit.

To ground these intra-chain advantages in the actual reward signal, we
normalize them to sum to the environment advantage:
\begin{align}
  \tilde{A}_{t,k}
    &= A_{t,k}^{\text{intra}} \cdot
       \frac{\hat{A}_t}
            {\sum_{j=1}^{K} A_{t,j}^{\text{intra}} + \delta},
  \label{eq:intra_norm}
\end{align}
where $\delta$ is a small constant for numerical stability.

The intra-chain value function is trained with the loss:
\begin{align}
  L^{\text{intra}}(\psi)
    &= \E_t\!\left[\bigl(
         V_\psi(\mathbf{s}_t, \mathbf{z}_{t,0}, 0)
       - V_\psi(\mathbf{s}_t, \mathbf{z}_{t,K}, K)
       - \hat{A}_t
       \bigr)^2\right],
  \label{eq:intra_loss}
\end{align}
which ensures the total value change across the chain predicts the
environment advantage.  The per-step PPO loss (\Cref{eq:perstep_clip})
then uses $\tilde{A}_{t,k}$ in place of $w_k \cdot \hat{A}_t$.

This approach extends temporal credit assignment \emph{inside} the
generative chain, analogous to how GAE assigns credit across
environment timesteps.  Unlike the weight-splitting schemes
(\Cref{sec:method:weights}), intra-chain credit is state- and
latent-dependent, grounded in learned value estimates, and need not
sum to a fixed proportion.

% ------------------------------------------------------------------
\subsection{Algorithm Summary}
\label{sec:method:algorithm}

\Cref{alg:perstep_ppo} summarizes the per-step PPO training loop.

\begin{algorithm}[t]
\caption{Per-Step PPO for Stochastic Flow Policies}
\label{alg:perstep_ppo}
\KwIn{Initial policy $\theta_0$, weight parameters (depending on scheme),
      clip parameter $\epsilon$, number of flow steps $K$, noise scale $\sigma$}
\For{iteration $= 1, 2, \ldots$}{
  Collect rollout batch using $\pi_{\theta_{\text{old}}}$\;
  \For{each transition $(\mathbf{s}_t, \mathbf{a}_t, \hat{A}_t)$}{
    Store intermediate latents $\{\mathbf{z}_{t,k}\}_{k=0}^{K}$\;
  }
  Compute GAE advantages $\{\hat{A}_t\}$\;
  \For{PPO epoch $= 1, \ldots, N_{\text{epochs}}$}{
    \For{minibatch of transitions}{
      \For{$k = K, \ldots, 1$}{
        Compute $r_{t,k}$ via \Cref{eq:perstep_ratio}\;
        Compute $w_k$ (or $w_{t,k}$) via chosen scheme\;
        $A_{t,k} \leftarrow w_k \cdot \hat{A}_t$\;
      }
      $L \leftarrow L^{\text{PS-CLIP}}$ \hfill (\Cref{eq:perstep_clip})\;
      Update $\theta$ (and $\phi$ or $\boldsymbol{\alpha}$ if applicable)
        via gradient ascent on $L$\;
    }
  }
}
\end{algorithm}

\section{Experimental Setup}
\label{sec:experiments}

% ------------------------------------------------------------------
\subsection{Environments}
\label{sec:experiments:envs}

We evaluate on two continuous-control locomotion tasks from the Gymnasium
MuJoCo suite~\citep{towers2024gymnasium}:
\begin{itemize}[nosep,leftmargin=*]
  \item \textbf{HalfCheetah-v5} ($d_s{=}17$, $d_a{=}6$): a planar running
        task with dense velocity reward.
  \item \textbf{Hopper-v5} ($d_s{=}11$, $d_a{=}3$): a single-legged hopping
        task requiring balance; episodes terminate on falls.
\end{itemize}
\todo{Optionally add Ant-v5 or AntMaze results.}

% ------------------------------------------------------------------
\subsection{Methods Compared}
\label{sec:experiments:methods}

We compare the following PPO variants, all using the same stochastic flow
policy architecture:
\begin{enumerate}[nosep,leftmargin=*]
  \item \textbf{Holistic PPO} (\Cref{sec:method:holistic}): standard PPO with
        a single ratio computed from the full chain log-probability.
  \item \textbf{Per-Step Uniform}: per-step PPO with uniform weights
        $w_k = 1/K$.
  \item \textbf{Per-Step Learned Global}: per-step PPO with learned global
        logits $\boldsymbol{\alpha}$.
  \item \textbf{Per-Step State-Dependent}: per-step PPO with the
        state-dependent weighting network $h_\phi$.
\end{enumerate}

% ------------------------------------------------------------------
\subsection{Architecture and Hyperparameters}
\label{sec:experiments:hyperparams}

The velocity field $v_\theta$ is parameterized as a 3-layer MLP with hidden
dimension~256 and SiLU activations.  The step index $k$ is embedded via a
sinusoidal positional encoding and concatenated with the state and latent
inputs.  The state-dependent weighting network $h_\phi$ is a 2-layer MLP with
hidden dimension~64.

Key hyperparameters are listed in \Cref{tab:hyperparams}.

\begin{table}[t]
  \centering
  \caption{Hyperparameters used across all experiments.}
  \label{tab:hyperparams}
  \begin{tabular}{@{}lc@{}}
    \toprule
    \textbf{Hyperparameter} & \textbf{Value} \\
    \midrule
    Flow steps $K$                  & 4 (default) \\
    Per-step noise $\sigma$         & 0.1 \\
    PPO clip $\epsilon$             & 0.2 \\
    GAE $\lambda$                   & 0.95 \\
    Discount $\gamma$               & 0.99 \\
    PPO epochs                      & 10 \\
    Minibatch size                  & 64 \\
    Rollout length                  & 2048 steps \\
    Learning rate (policy)          & $3 \times 10^{-4}$ \\
    Learning rate (value)           & $1 \times 10^{-3}$ \\
    Learning rate (weights $\alpha/\phi$) & $1 \times 10^{-3}$ \\
    Entropy coefficient             & 0.0 \\
    Value loss coefficient          & 0.5 \\
    Max gradient norm               & 0.5 \\
    Total environment steps         & $1 \times 10^{6}$ \\
    \bottomrule
  \end{tabular}
\end{table}

% ------------------------------------------------------------------
\subsection{Evaluation Protocol}
\label{sec:experiments:eval}

We train each method with 3 random seeds and evaluate every \todo{N} training
steps by running 10 deterministic episodes (using the mean velocity field
without per-step noise, i.e., $\sigma{=}0$ at evaluation time).  We report the
mean and standard deviation of the episodic return across seeds.

% ------------------------------------------------------------------
\subsection{Ablations}
\label{sec:experiments:ablations}

We conduct the following ablation studies:
\begin{itemize}[nosep,leftmargin=*]
  \item \textbf{Number of flow steps $K$}: we vary $K \in \{2, 4, 8\}$ to
        assess how the gap between holistic and per-step PPO changes with chain
        length.
  \item \textbf{Weighting mode}: we compare uniform, learned-global, and
        state-dependent weighting.
  \item \textbf{Noise scale $\sigma$}: we vary
        $\sigma \in \{0.01, 0.1, 0.5\}$ to test sensitivity to per-step
        stochasticity.
\end{itemize}

\section{Results}
\label{sec:results}

% ------------------------------------------------------------------
\subsection{Main Comparison}
\label{sec:results:main}

\Cref{tab:main_results} presents the final evaluation returns for all methods.
\todo{Fill in with actual numbers.}

\begin{table}[t]
  \centering
  \caption{Mean evaluation return ($\pm$ std) over 3 seeds after $1\text{M}$
           environment steps.  Best result per environment in \textbf{bold}.}
  \label{tab:main_results}
  \begin{tabular}{@{}lcc@{}}
    \toprule
    \textbf{Method} & \textbf{HalfCheetah-v5} & \textbf{Hopper-v5} \\
    \midrule
    Holistic PPO              & \todo{XXX $\pm$ YYY} & \todo{XXX $\pm$ YYY} \\
    Per-Step Uniform          & \todo{XXX $\pm$ YYY} & \todo{XXX $\pm$ YYY} \\
    Per-Step Learned Global   & \todo{XXX $\pm$ YYY} & \todo{XXX $\pm$ YYY} \\
    Per-Step State-Dependent  & \todo{XXX $\pm$ YYY} & \todo{XXX $\pm$ YYY} \\
    \bottomrule
  \end{tabular}
\end{table}

\todo{Discuss which methods outperform the baseline and by how much.  Note any
environment-specific differences.}

% ------------------------------------------------------------------
\subsection{Training Curves}
\label{sec:results:curves}

\Cref{fig:training_curves} shows the evaluation return over the course of
training.

\begin{figure}[t]
  \centering
  \begin{subfigure}[t]{0.48\textwidth}
    \centering
    \includegraphics[width=\textwidth]{figs/halfcheetah_training_curves.pdf}
    \caption{HalfCheetah-v5}
    \label{fig:curve_halfcheetah}
  \end{subfigure}
  \hfill
  \begin{subfigure}[t]{0.48\textwidth}
    \centering
    \includegraphics[width=\textwidth]{figs/hopper_training_curves.pdf}
    \caption{Hopper-v5}
    \label{fig:curve_hopper}
  \end{subfigure}
  \caption{Evaluation return vs.\ environment steps for all methods.  Solid
           lines show the mean over 3 seeds; shaded regions show $\pm 1$ std.}
  \label{fig:training_curves}
\end{figure}

\todo{Describe learning speed differences.  Does per-step PPO learn faster or
more stably?}

% ------------------------------------------------------------------
\subsection{Weight Profile Analysis}
\label{sec:results:weights}

For the learned-global weighting scheme, \Cref{fig:weight_profile} visualizes
the learned weight distribution $w_k$ across the $K$ flow steps after
training.

\begin{figure}[t]
  \centering
  \includegraphics[width=0.5\textwidth]{figs/weight_profile.pdf}
  \caption{Learned global weights $w_k$ (softmax of $\boldsymbol{\alpha}$)
           across flow steps for HalfCheetah-v5.  \todo{Describe pattern.}}
  \label{fig:weight_profile}
\end{figure}

\todo{Discuss which steps receive the most credit.  Do early steps (close to
the final action) or late steps (close to noise) get higher weight?  Is the
pattern consistent across environments?}

% ------------------------------------------------------------------
\subsection{Per-Step Clip Fraction Analysis}
\label{sec:results:clipfrac}

\Cref{fig:clip_fractions} compares the per-step clip fractions (fraction of
samples where $r_{t,k}$ is clipped) between holistic and per-step PPO.

\begin{figure}[t]
  \centering
  \includegraphics[width=0.5\textwidth]{figs/clip_fractions.pdf}
  \caption{Per-step clip fractions during training.  \todo{Describe the
           difference between holistic and per-step clipping.}}
  \label{fig:clip_fractions}
\end{figure}

\todo{Discuss whether per-step clipping leads to more uniform and moderate clip
fractions compared to holistic clipping.}

% ------------------------------------------------------------------
\subsection{Ablation: Number of Flow Steps $K$}
\label{sec:results:ablation_K}

\Cref{tab:ablation_K} shows results as we vary the number of flow steps.

\begin{table}[t]
  \centering
  \caption{Evaluation return on HalfCheetah-v5 for different $K$.}
  \label{tab:ablation_K}
  \begin{tabular}{@{}lccc@{}}
    \toprule
    \textbf{Method} & $K{=}2$ & $K{=}4$ & $K{=}8$ \\
    \midrule
    Holistic PPO              & \todo{---} & \todo{---} & \todo{---} \\
    Per-Step Uniform          & \todo{---} & \todo{---} & \todo{---} \\
    Per-Step Learned Global   & \todo{---} & \todo{---} & \todo{---} \\
    \bottomrule
  \end{tabular}
\end{table}

\todo{Discuss how the gap between holistic and per-step changes with $K$.
Hypothesis: per-step PPO should benefit more for larger $K$, where the holistic
ratio becomes more extreme.}

% ------------------------------------------------------------------
\subsection{Ablation: Noise Scale $\sigma$}
\label{sec:results:ablation_sigma}

\begin{table}[t]
  \centering
  \caption{Evaluation return on HalfCheetah-v5 for different noise scales
           $\sigma$, using per-step uniform PPO with $K{=}4$.}
  \label{tab:ablation_sigma}
  \begin{tabular}{@{}lccc@{}}
    \toprule
    \textbf{Method} & $\sigma{=}0.01$ & $\sigma{=}0.1$ & $\sigma{=}0.5$ \\
    \midrule
    Holistic PPO     & \todo{---} & \todo{---} & \todo{---} \\
    Per-Step Uniform & \todo{---} & \todo{---} & \todo{---} \\
    \bottomrule
  \end{tabular}
\end{table}

\todo{Discuss sensitivity to $\sigma$.  Smaller $\sigma$ concentrates density
and may produce sharper ratios; larger $\sigma$ makes per-step densities
broader.}

% ------------------------------------------------------------------
\subsection{Discussion}
\label{sec:results:discussion}

\todo{Synthesize findings.  When does per-step PPO help most?  Is
state-dependent weighting worth the additional complexity?  What are failure
modes?}

\section{Conclusion}
\label{sec:conclusion}

\paragraph{Summary.}
We studied the question of how PPO should interact with the multi-step
stochastic process inside a flow-based generative policy.  We formalized
per-step PPO, which applies clipped importance-sampling ratios independently at
each generative step, and introduced weighted credit-assignment schemes---uniform,
learned-global, and state-dependent---to modulate gradient signal across steps.
\todo{State the main empirical finding in 1--2 sentences.  E.g., ``Per-step PPO
with learned global weights consistently outperformed holistic PPO on
HalfCheetah-v5 and Hopper-v5, with the largest gains observed for $K{=}8$
flow steps.''}

\paragraph{Limitations.}
Our study has several limitations:
\begin{itemize}[nosep,leftmargin=*]
  \item \textbf{Computational overhead.}  Per-step PPO requires storing
        intermediate latents $\{\mathbf{z}_{t,k}\}$ for every transition and
        computing $K$ ratios per sample, increasing memory and compute relative
        to holistic PPO.
  \item \textbf{Limited environments.}  We evaluated on two MuJoCo locomotion
        tasks.  Results may differ on higher-dimensional or sparse-reward tasks
        such as AntMaze or dexterous manipulation.
  \item \textbf{Fixed noise schedule.}  We used a constant $\sigma$ across all
        steps; an adaptive or learned noise schedule could interact differently
        with per-step clipping.
  \item \textbf{No theoretical guarantees.}  We lack a formal analysis of when
        per-step clipping provides better trust-region properties than holistic
        clipping.
\end{itemize}

\paragraph{Future Work.}
Several directions are promising:
\begin{itemize}[nosep,leftmargin=*]
  \item \textbf{Learned noise schedule.}  Jointly learning $\sigma_k$ per step
        could improve both policy expressiveness and the quality of per-step
        ratios.
  \item \textbf{Scaling to larger $K$.}  With many flow steps ($K \geq 16$),
        the advantage of per-step clipping should grow; evaluating this regime
        is important.
  \item \textbf{Broader benchmarks.}  Testing on image-based manipulation,
        multi-task settings, and sparse-reward environments would strengthen
        generality claims.
  \item \textbf{Theoretical analysis.}  Characterizing the bias-variance
        trade-off of per-step vs.\ holistic clipping, and conditions under
        which weighted credit assignment provably improves convergence, would
        complement our empirical results.
\end{itemize}


% ---------- references ----------
\bibliographystyle{plainnat}
\bibliography{references}

\end{document}
